\documentclass[10pt]{article}
\usepackage[T1]{fontenc}
\usepackage[utf8]{inputenc}
\usepackage[french]{babel}
\usepackage{amsmath,amssymb,mathrsfs,mathpazo}
\usepackage{geometry}
\geometry{a4paper,margin=2cm}
\usepackage{graphicx,array,tabularx,tabularray,colortbl,multirow,multicol,booktabs,enumitem,pifont,ulem,fancybox,diagbox,cancel}
\usepackage{pstricks,pstricks-add,pst-plot,pst-tree,pst-node,pst-text,pst-func,pst-math,pst-bspline,pst-3dplot,pst-eucl,pst-solides3d}
\let\pscalculate\undefined
\usepackage{pst-bezier}
\usepackage{tikz,pgfplots,tkz-tab}
\usetikzlibrary{arrows,arrows.meta,patterns,calc,positioning,trees,decorations.markings,shapes.geometric,shapes.misc,shapes.arrows,intersections,decorations.pathmorphing,decorations.pathreplacing,backgrounds,shadows,babel}
\pgfplotsset{compat=1.18}
\usepgfplotslibrary{fillbetween}
\usepackage[np]{numprint}
\usepackage{listings,makecell,pas-tableur,twoopt}
\let\barre\undefined
\usepackage[pstricks]{tablvar}
\definecolor{darkgreen}{RGB}{0,100,0}
\providecommand{\red}{\color{red}}
\providecommand{\blue}{\color{blue}}
\providecommand{\green}{\color{green}}
\providecommand{\black}{\color{black}}
\providecommand{\white}{\color{white}}
\providecommand{\gray}{\color{gray}}
\providecommand{\cyan}{\color{cyan}}
\providecommand{\magenta}{\color{magenta}}
\providecommand{\yellow}{\color{yellow}}
\providecommand{\violet}{\color{violet}}
\providecommand{\orange}{\color{orange}}
\providecommand{\euro}{\texteuro}
\providecommand{\EUR}{\texteuro}
\providecommand{\ieme}{\textsuperscript{e}}
\providecommand{\degres}{\ensuremath{^\circ}}
\providecommand{\degree}{\ensuremath{^\circ}}
\providecommand{\vv}[1]{\overrightarrow{#1}}
\providecommand{\Coord}[1]{#1}
\providecommand{\decosix}{}
\providecommand{\starredbullet}{\ensuremath{\bullet}}
\pagestyle{empty}
\setlength{\parindent}{0pt}
\usepackage{adjustbox}\DeclareUnicodeCharacter{2236}{:}\DeclareUnicodeCharacter{2190}{\ensuremath{\leftarrow}}
\providecommand{\R}{\mathbb{R}}
\providecommand{\N}{\mathbb{N}}
\providecommand{\Z}{\mathbb{Z}}
\providecommand{\Q}{\mathbb{Q}}
\providecommand{\D}{\mathbb{D}}
\providecommand{\C}{\mathbb{C}}
\providecommand{\e}{{\mathrm{e}}}
\providecommand{\vect}[1]{\overrightarrow{#1}}
\providecommand{\vectt}[1]{\overrightarrow{\mathrm{#1}}}
\providecommand{\vvt}[1]{\overrightarrow{\mathrm{#1}}}
\providecommand{\np}[2][]{{#2~#1}}
\providecommand{\vv}[1]{{\overrightarrow{#1}}}
\renewcommand{\d}{{\mathrm{d}}}
\providecommand{\degres}{{^\circ}}
\providecommand{\degre}{{^\circ}}
\providecommand{\blue}{\color{blue}}
\providecommand{\red}{\color{red}}
\providecommand{\enskip}{\quad}
\providecommand{\llbracket}{\lbrack\!\lbrack}
\providecommand{\rrbracket}{\rbrack\!\rbrack}
\providecommand{\euro}{{\text{€}}}
\providecommand{\no}{{\text{n°}}}
\providecommand{\interval}[2]{{[#1;#2]}}
\providecommand{\intervFO}[2]{{[#1;#2[}}
\providecommand{\dotfill}{\ldots}

\providecommand{\ds}{\displaystyle}
\providecommand{\textding}[1]{#1}
\providecommand{\Oij}{\(\mathrm{O};\vec{\imath},\vec{\jmath}\)}
\providecommand{\Oijk}{\(\mathrm{O};\vec{\imath},\vec{\jmath},\vec{k}\)}
\providecommand{\Ouv}{\(\mathrm{O};\vec{u},\vec{v}\)}
\providecommand{\intervFF}[2]{\ensuremath{[#1;#2]}}
\renewcommand{\euro}{\ifmmode\text{\texteuro}\else\texteuro\fi}
\DeclareUnicodeCharacter{00B0}{\ifmmode^{\circ}\else\textdegree\fi}
\DeclareUnicodeCharacter{20AC}{\ifmmode\text{\texteuro}\else\texteuro\fi}
% Énoncés : Éducation nationale. Transcription et corrigés : APMEP, Denis Vergès et contributeurs. Découpage et classement : CoopMaths, Sébastien Lozano, Éric Elter et contributeurs. Adaptation de la présentation : Maths Explorer.
% CC BY-SA : https://coopmaths.fr/www/about/
% https://coopmaths.fr/alea/static/eam/2026/tex/eam_technologique_2026_09_antilles_1_cor.tex
\begin{document}
\noindent\textit{APMEP / CoopMaths. CC BY-SA. Présentation adaptée par Maths Explorer.}\par\medskip
\subsection*{Partie A}

\begin{enumerate}
	\item $\mathcal{C}_R$ est une droite passant par l'origine : $R$ est donc une fonction linéaire, de la forme $R(x) = ax$.

	D'après les deux points de coordonnées $(20~;~\np{6000})$ et $(40~;~\np{12000})$ marqués sur le graphique :
	\[a = \dfrac{\np{12000} - \np{6000}}{40 - 20} = \dfrac{\np{6000}}{20} = 300.\]
	
Pour tout $x$ de $[0~;~40]$ : $R(x) = 300x$.

	\item L'entreprise réalise un bénéfice lorsque la recette dépasse le coût, c'est-à-dire lorsque
	$\mathcal{C}_R$ est au-dessus de $\mathcal{C}_C$.

	Par lecture graphique, $\mathcal{C}_R$ est au-dessus de $\mathcal{C}_C$ entre les deux
	points d'intersection des courbes, d'abscisses environ 18 et 39.
	L'entreprise réalise donc un bénéfice pour $x$ compris entre environ
	18 et 39 tonnes.

	\item Le bénéfice est maximal lorsque l'écart vertical
	entre $\mathcal{C}_R$ (au-dessus) et $\mathcal{C}_C$ (en dessous) est le plus grand.

	Avec la précision permise par le graphique, cet écart est maximal pour
	$x \approx 30$ ; le bénéfice maximal est alors d'environ \np{3400}~\euro.
\end{enumerate}

\subsection*{Partie B}

\emph{Pour $x \in [0~;~40]$, $B(x) = -x^3+54x^2-540x-\np{2000}$.}

\begin{enumerate}
	\item Pour tout réel $x$ de $[0~;~40]$ :
	\[B'(x) = -3x^2 + 54 \times 2x - 540 = -3x^2 + 108x - 540.\]
	D'autre part :
	\[(-3x+18)(x-30) = -3x^2 - 90x + 18x - 540 = -3x^2 + 108x - 540.\]
	Donc, pour tout $x$ de $[0 ; 40]$, $B'(x) = (-3x+18)(x-30)$.

	\item $-3x + 18 = 0 \iff x = 6$ (coefficient $-3 < 0$ : positif avant, négatif après) ;

	$x - 30 = 0 \iff x = 30$ (coefficient $1 > 0$ : négatif avant, positif après).

	\begin{center}
		\begin{tikzpicture}
			\tkzTabInit[lgt=3.4,espcl=2.6]
			{$x$/1,$-3x+18$/1,$x-30$/1,$B'(x)$/1,$B$/2}{$0$,$6$,$30$,$40$}
			\tkzTabLine{,+,z,-,,-,}
			\tkzTabLine{,-,,,z,+,}
			\tkzTabLine{,-,z,+,z,-,}
			\tkzTabVar{+/{$-\np{2000}$},-/{$-\np{3512}$},+/{$\np{3400}$},-/{$-\np{1200}$}}
		\end{tikzpicture}
	\end{center}

La fonction $B$ est décroissante sur $[0~;~6]$, croissante sur
	$[6~;~30]$, puis décroissante sur $[30~;~40]$.

	\item D'après le tableau de variations, le bénéfice est maximal pour $x = 30$ tonnes de compote produites et vendues. Il vaut alors $B(30) = \np{3400}$~(\euro).
\end{enumerate}

\medskip

\end{document}
